\section{Metodología}

Se aplicó un procedimiento descriptivo y reproducible. El análisis consideró 3.000 incidentes registrados durante 2025 y los complementó con eventos y encuestas simuladas para calcular las métricas de calidad en uso.

El trabajo se realizó en cuatro pasos: análisis y clasificación de incidentes; definición de tareas, atributos y métricas; generación y procesamiento automatizado de evidencia; e interpretación de KPI y formulación del plan de mejora. Python se utilizó para clasificación, simulación y cálculo; React y Flutter para representar las tareas de los usuarios; y los servicios locales para registrar los fallos controlados.

\begin{table}[H]
\centering
\caption{Fuentes empleadas en la evaluación}
\begin{tabularx}{\textwidth}{p{3cm}p{2.2cm}YY}
\toprule
Fuente & Cantidad & Tipo de dato & Uso \\
\midrule
Incidentes & 3.000 & Registro de casos & Clasificación ISO/IEC 25022 \\
Eventos & 6.000 & Simulación, semilla 25022 & Efectividad, eficiencia, riesgo y contexto \\
Encuestas & 2.000 & Simulación, semilla 25022 & Satisfacción y abandono \\
Catálogo & 10 métricas & Diseño del estudio & Fórmulas, metas y semáforo \\
\bottomrule
\end{tabularx}
\end{table}

Los incidentes se normalizaron y clasificaron según su característica principal de calidad en uso. Los eventos y encuestas se procesaron por separado para evitar que los datos simulados alteraran el registro de casos y para mantener comparables los resultados de cada métrica.
